\subsection{مبانی نظری و فرمول‌بندی ریاضی}
فرض کنید $\mathbf{x} \in \mathbb{R}^{d_{\text{in}}}$ بردار فعال‌سازی ورودی به لایه باشد. تقریب قطری اطلاعات فیشر فعال‌سازی از طریق گشتاور دوم نمونه‌های کالیبراسیون محاسبه می‌شود:
\begin{equation}
\mathbf{F}_i = \mathbb{E}\left[\mathbf{x}_i^2\right] = \frac{1}{B} \sum_{k=1}^B (\mathbf{X}_{k, i})^2, \quad i \in \{1, \dots, d_{\text{in}}\}
\end{equation}
با تعریف ماتریس قطری $\mathbf{D}_F = \text{diag}\left(\sqrt{\mathbf{F}_1}, \dots, \sqrt{\mathbf{F}_{d_{\text{in}}}}\right)$ با اعمال کران پایینی عددی $\epsilon \ge 10^{-8}$ جهت پایداری، ماتریس وزن وزن‌دهی می‌شود:
\begin{equation}
\mathbf{W}_{\text{weighted}} = \mathbf{W} \mathbf{D}_F \in \mathbb{R}^{d_{\text{out}} \times d_{\text{in}}}
\end{equation}
سپس تجزیه مقادیر منفرد بر روی ماتریس وزن‌دهی‌شده اجرا می‌شود:
\begin{equation}
\mathbf{W}_{\text{weighted}} = \mathbf{U} \mathbf{\Sigma} \mathbf{V}^\top
\end{equation}
با انتخاب رتبه فشرده‌سازی $r \le \min(d_{\text{in}}, d_{\text{out}})$، مقادیر منفرد به فاکتورها تقسیم شده و فاکتور راست با معکوس فیشر بی‌وزن می‌گردد:
\begin{align}
\mathbf{U}_r &= \mathbf{U}_{:, :r} \mathbf{\Sigma}_{:r}^{1/2} \in \mathbb{R}^{d_{\text{out}} \times r} \\
\mathbf{V}_r &= \mathbf{\Sigma}_{:r}^{1/2} \mathbf{V}_{:, :r}^\top \mathbf{D}_F^{-1} \in \mathbb{R}^{r \times d_{\text{in}}}
\end{align}
بازسازی حاصل به فرم $\mathbf{W} \approx \mathbf{U}_r \mathbf{V}_r$ تضمین می‌کند که خطای بازسازی فروبنیوس در جهات دارای اطلاعات فیشر بالا به حداقل ممکن کاهش یابد:
\begin{equation}
\mathcal{E} = \frac{\|\mathbf{W} - \mathbf{U}_r \mathbf{V}_r\|_F}{\|\mathbf{W}\|_F}
\end{equation}
محاسبه خروجی لایه فشرده در زمان استنتاج به دو ضرب ماتریسی سبک با مرتبه پارامتری $r \cdot (d_{\text{in}} + d_{\text{out}})$ تبدیل می‌شود:
\begin{equation}
\mathbf{y} = (\mathbf{x} \mathbf{V}_r^\top) \mathbf{U}_r^\top + \mathbf{b}
\end{equation}

که در آن $\sigma_i(\mathbf{W})$ بیانگر $i$-امین مقدار منفرد مرتب‌شده ماتریس به صورت نزولی است. نسبت حفظ انرژی طیفی $\eta_r = \frac{\sum_{i=1}^r \sigma_i^2}{\sum_{i=1}^{d} \sigma_i^2}$ تضمین می‌کند که با انتخاب آستانه $\tau \ge 0.90$، بیش از نود درصد از پراکندگی واریانس داده‌ها حفظ شده و افت سیگنال در انتشار گرادیان رو به عقب به حداقل ممکن کاهش می‌یابد.

\paragraph{تحلیل مجانبی و پیچیدگی اشتراک حافظه سخت‌افزاری}
از منظر پیچیدگی محاسباتی و انتقال داده‌ها در سلسومراتب حافظه (حافظه سراسری \lr{HBM} به حافظه محلی \lr{SRAM})، عملیات ضرب ماتریسی در این مکانیزم نیازمند تخصیص بافر موقت با ابعاد فاکتورگیری‌شده است. برای پردازش دنباله‌ای به طول $N$ با بعد مخفی $D$، بار محاسباتی هسته از فرم درجه دو به رابطه زیر تقلیل می‌یابد:
\begin{equation}
\label{eq:complexity_fwsvd_compressor}
\mathcal{C}_{\text{fwsvd\_compressor}} = \mathcal{O}\left( B \cdot N \cdot D \cdot r + B \cdot H \cdot N \cdot d_k \right) \ll \mathcal{O}\left( B \cdot H \cdot N^2 \cdot d_k \right)
\end{equation}
این بهبود مرتبه پیچیدگی، با کاهش چشمگیر حجم تبادل تانسورها در گذرگاه‌های ارتباطی کارت گرافیک، مانع از اشباع پهنای باند حافظه (پدیده \lr{Memory-Bound}) در فرآیند استنتاج طولانی‌مدت می‌گردد.
